\chapter{Triângulos e quadriláteros}\label{ch:g6:shapes}

Triângulos e \omterm{def:g4:geometry:polygon}{quadriláteros} são as formas básicas da geometria. O
objetivo deste capítulo é conhecer as famílias especiais deles ---
\omterm{def:g6:shapes:triangles}{isósceles}, \omterm{def:g6:shapes:triangles}{equiláteros}, retângulos; retângulos, \omterm{def:g6:shapes:quadrilaterals}{losangos}, quadrados ---
pelas suas \emph{propriedades definidoras}, e construí-los com precisão
com régua e compasso.

\section{Triângulos}

\begin{definition}[Triângulos especiais]\label{def:g6:shapes:triangles}
Um triângulo $ABC$ tem três \omterm{def:g5:solids:def}{vértices} e três lados. Ele é:
\begin{itemize}
  \item \emph{isósceles}\index{triângulo isósceles} em $A$ quando
        $AB = AC$ (dois lados iguais);
  \item \emph{equilátero}\index{triângulo equilátero} quando
        $AB = BC = CA$ (três lados iguais);
  \item \emph{retângulo}\index{triângulo retângulo} em $A$ quando o
        ângulo $\widehat{BAC}$ é reto. O lado $[BC]$, oposto ao \omterm{def:g3:shapes:rightangle}{ângulo
        reto}, chama-se \emph{hipotenusa}.
\end{itemize}
\end{definition}

\begin{omfigure}
\begin{tikzpicture}[scale=0.85]
  % isósceles
  \draw[thick] (0,0) -- (2.4,0) -- (1.2,2.1) -- cycle;
  \draw (0.55,0.96) -- (0.72,0.86);
  \draw (1.85,0.96) -- (1.68,0.86);
  \node[below, font=\small] at (1.2,-0.15) {isósceles};
  \node[above, font=\small] at (1.2,2.1) {$A$};
  \node[below left, font=\small] at (0,0) {$B$};
  \node[below right, font=\small] at (2.4,0) {$C$};
  % equilátero
  \begin{scope}[xshift=4.6cm]
  \draw[thick] (0,0) -- (2.4,0) -- (1.2,2.08) -- cycle;
  \draw (1.13,0.05) -- (1.27,-0.05);
  \draw (0.55,0.99) -- (0.72,0.89);
  \draw (1.85,0.99) -- (1.68,0.89);
  \node[below, font=\small] at (1.2,-0.25) {equilátero};
  \end{scope}
  % retângulo
  \begin{scope}[xshift=9.2cm]
  \draw[thick] (0,0) -- (2.6,0) -- (0,1.8) -- cycle;
  \draw[thin] (0.3,0) -- (0.3,0.3) -- (0,0.3);
  \node[below, font=\small] at (1.3,-0.15) {retângulo};
  \node[omThm, font=\small, rotate=-34] at (1.5,1.05) {hipotenusa};
  \end{scope}
\end{tikzpicture}

{\small Os três triângulos especiais. Os tracinhos indicam lados
iguais; o quadradinho marca o \omterm{def:g3:shapes:rightangle}{ângulo reto}.}
\end{omfigure}

\begin{method}[Construir um triângulo a partir dos três lados]\label{met:g6:shapes:construct}
Para construir $ABC$ com $AB = 5$ cm, $AC = 4$ cm, $BC = 3$ cm:
\begin{enumerate}
  \item trace o \omterm{def:g6:lines:objects}{segmento} $[AB]$ de $5$ cm com a régua;
  \item trace um arco do círculo de centro $A$ e \omterm{def:g3:shapes:shapes}{raio} $4$ cm (todos os
        pontos a $4$ cm de $A$);
  \item trace um arco do círculo de centro $B$ e \omterm{def:g3:shapes:shapes}{raio} $3$ cm;
  \item o ponto $C$ é onde os dois arcos se cruzam; ligue-o a $A$ e a
        $B$.
\end{enumerate}
\end{method}

\begin{omfigure}
\begin{tikzpicture}[scale=1.0]
  \coordinate (A) at (0,0);
  \coordinate (B) at (5,0);
  \coordinate (C) at (3.2,2.4);
  \draw[thick] (A) -- (B);
  \draw[thick, omDef] (A) -- (C);
  \draw[thick, omThm] (B) -- (C);
  \draw[omDef, thin] (C) arc[start angle=36.9, delta angle=22, radius=4];
  \draw[omDef, thin] (C) arc[start angle=36.9, delta angle=-22, radius=4];
  \draw[omThm, thin] (C) arc[start angle=126.9, delta angle=22, radius=3];
  \draw[omThm, thin] (C) arc[start angle=126.9, delta angle=-22, radius=3];
  \fill (A) circle (1.5pt) node[below left, font=\small] {$A$};
  \fill (B) circle (1.5pt) node[below right, font=\small] {$B$};
  \fill (C) circle (1.5pt) node[above, font=\small] {$C$};
  \node[omDef, font=\small] at (1.2,1.5) {$4$ cm};
  \node[omThm, font=\small] at (4.6,1.4) {$3$ cm};
  \node[below, font=\small] at (2.5,0) {$5$ cm};
\end{tikzpicture}

{\small Construção de um triângulo com o compasso: o \omterm{def:g5:solids:def}{vértice} $C$ é o
ponto de cruzamento de dois arcos. (No \cref{ch:g7:triangles} você vai
ver quando essa construção é possível.)}
\end{omfigure}

\section{Quadriláteros}

\begin{definition}[Quadriláteros especiais]\label{def:g6:shapes:quadrilaterals}
Um \omterm{def:g4:geometry:polygon}{quadrilátero} tem quatro \omterm{def:g5:solids:def}{vértices} e quatro lados. Ele é um:
\begin{itemize}
  \item \emph{retângulo}\index{retângulo} quando os quatro ângulos dele
        são retos;
  \item \emph{losango}\index{losango} quando os quatro lados dele têm o
        mesmo comprimento;
  \item \emph{quadrado}\index{quadrado} quando é as duas coisas: quatro
        \omterm{def:g3:shapes:rightangle}{ângulos retos} \emph{e} quatro lados iguais.
\end{itemize}
\end{definition}

\begin{proposition}[Propriedades das diagonais]\label{prop:g6:shapes:diagonals}
Nesses \omterm{def:g4:geometry:polygon}{quadriláteros}, as duas diagonais se cortam no ponto médio comum
delas e, além disso:
\begin{itemize}
  \item num retângulo, as diagonais têm o \emph{mesmo comprimento};
  \item num \omterm{def:g6:shapes:quadrilaterals}{losango}, as diagonais são \emph{\omterm{def:g6:lines:perp}{perpendiculares}};
  \item num quadrado, as duas coisas ao mesmo tempo.
\end{itemize}
\end{proposition}
\admitted

\begin{omfigure}
\begin{tikzpicture}[scale=0.8]
  % retângulo, marcas de ângulo reto voltadas para dentro
  \draw[thick] (0,0) rectangle (3,1.9);
  \draw[omThm] (0,0) -- (3,1.9);
  \draw[omThm] (0,1.9) -- (3,0);
  \foreach \x/\y/\dx/\dy in {0/0/1/1, 3/0/-1/1, 0/1.9/1/-1, 3/1.9/-1/-1}
    \draw[thin] (\x+0.22*\dx,\y) -- (\x+0.22*\dx,\y+0.22*\dy)
      -- (\x,\y+0.22*\dy);
  \node[below, font=\small] at (1.5,-0.2) {retângulo};
  % losango
  \begin{scope}[xshift=5.1cm, yshift=0.95cm]
  \draw[thick] (0,0) -- (1.6,1.0) -- (3.2,0) -- (1.6,-1.0) -- cycle;
  \draw[omThm] (0,0) -- (3.2,0);
  \draw[omThm] (1.6,1.0) -- (1.6,-1.0);
  \draw[thin] (1.6,0) ++(0.18,0) -- ++(0,0.18) -- ++(-0.18,0);
  \node[below, font=\small] at (1.6,-1.2) {losango};
  \end{scope}
  % quadrado, com as marcas de ângulo reto também
  \begin{scope}[xshift=9.8cm]
  \draw[thick] (0,0) rectangle (1.9,1.9);
  \draw[omThm] (0,0) -- (1.9,1.9);
  \draw[omThm] (0,1.9) -- (1.9,0);
  \foreach \x/\y/\dx/\dy in {0/0/1/1, 1.9/0/-1/1, 0/1.9/1/-1,
                             1.9/1.9/-1/-1}
    \draw[thin] (\x+0.22*\dx,\y) -- (\x+0.22*\dx,\y+0.22*\dy)
      -- (\x,\y+0.22*\dy);
  \node[below, font=\small] at (0.95,-0.2) {quadrado};
  \end{scope}
\end{tikzpicture}

{\small Os três \omterm{def:g4:geometry:polygon}{quadriláteros} especiais com as suas diagonais (em
vermelho). O quadrado herda as propriedades do retângulo e do
\omterm{def:g6:shapes:quadrilaterals}{losango}.}
\end{omfigure}

\begin{remark}
Um quadrado é automaticamente um retângulo (tem quatro \omterm{def:g3:shapes:rightangle}{ângulos retos}) e
automaticamente um \omterm{def:g6:shapes:quadrilaterals}{losango} (quatro lados iguais). Então ``retângulo''
não quer dizer ``que não é quadrado''! Em matemática, a forma mais
especial pertence à família mais geral.
\end{remark}

\begin{example}[Reconhecer uma forma pelas propriedades]\label{ex:g6:shapes:recognize}
Um \omterm{def:g4:geometry:polygon}{quadrilátero} tem quatro lados de comprimento $4$ cm e um \omterm{def:g3:shapes:rightangle}{ângulo
reto}. O que é ele? Quatro lados iguais fazem dele um \omterm{def:g6:shapes:quadrilaterals}{losango}. Num
\omterm{def:g6:shapes:quadrilaterals}{losango}, os ângulos opostos são iguais e os quatro ângulos somam
$360^\circ$; um \omterm{def:g3:shapes:rightangle}{ângulo reto} obriga então os quatro a serem retos (isso
será demonstrado com os paralelogramos no \cref{ch:g7:central}). Logo a
forma é um quadrado.
\end{example}

\begin{method}[Programas de construção]\label{met:g6:shapes:program}
Um \emph{programa de construção} é uma lista de instruções que qualquer
pessoa pode seguir para reproduzir uma figura exatamente. Escreva uma
instrução por linha, dê nome a cada ponto no momento em que ele é
criado, e use apenas: ``trace o \omterm{def:g6:lines:objects}{segmento} / a reta / o círculo\dots'',
``marque o ponto onde \dots\ se cruzam''. Teste o seu programa
seguindo-o ao pé da letra, palavra por palavra.
\end{method}

\begin{example}\label{ex:g6:shapes:program}
Programa para um quadrado $ABCD$ de lado $4$ cm:
\begin{enumerate}
  \item Trace um \omterm{def:g6:lines:objects}{segmento} $[AB]$ com $AB = 4$ cm.
  \item Trace a perpendicular a $(AB)$ em $A$; sobre ela, marque $D$ a
        $4$ cm de $A$ (escolha um lado da reta e mantenha-o).
  \item Trace a perpendicular a $(AB)$ em $B$; sobre ela, marque $C$ a
        $4$ cm de $B$, do mesmo lado que $D$.
  \item Ligue $C$ e $D$.
\end{enumerate}
Confira com a régua: $DC$ também deve medir $4$ cm.
\end{example}

\section{Exercícios}

\begin{exercise}[$\star$]\label{exo:g6:shapes:1}
Construa um triângulo $ABC$ com $AB = 6$ cm, $AC = 5$ cm e
$BC = 4$ cm, seguindo o \cref{met:g6:shapes:construct}.
\end{exercise}

\begin{exercise}[$\star$]\label{exo:g6:shapes:2}
Construa um \omterm{def:g6:shapes:triangles}{triângulo isósceles} $DEF$ com $DE = DF = 5$ cm e
$EF = 3$ cm. Que dois ângulos desse triângulo parecem iguais? Meça-os
para conferir.
\end{exercise}

\begin{exercise}[$\star$]\label{exo:g6:shapes:3}
Construa um \omterm{def:g6:shapes:triangles}{triângulo equilátero} de lado $4$ cm usando apenas régua e
compasso. Meça um dos ângulos dele.
\end{exercise}

\begin{exercise}[$\star$]\label{exo:g6:shapes:4}
Classifique cada triângulo a partir dos dados (\omterm{def:g6:shapes:triangles}{isósceles}, \omterm{def:g6:shapes:triangles}{equilátero},
retângulo --- várias respostas podem valer):
(a) lados $7$, $7$, $7$; \
(b) lados $5$, $5$, $8$; \
(c) ângulos $90^\circ$, $45^\circ$, $45^\circ$; \
(d) ângulos $60^\circ$, $60^\circ$, $60^\circ$.
\end{exercise}

\begin{exercise}[$\star$]\label{exo:g6:shapes:5}
Num \omterm{def:g6:shapes:triangles}{triângulo retângulo}, qual lado é a hipotenusa? Trace um triângulo
$MNP$ retângulo em $N$ e nomeie a hipotenusa dele. Que \omterm{def:g5:solids:def}{vértice} a
hipotenusa nunca toca?
\end{exercise}

\begin{exercise}[$\star$]\label{exo:g6:shapes:6}
Verdadeiro ou falso? Justifique cada resposta em uma frase.
\begin{enumerate}
  \item Todo quadrado é um retângulo.
  \item Todo retângulo é um quadrado.
  \item Todo quadrado é um \omterm{def:g6:shapes:quadrilaterals}{losango}.
  \item Um \omterm{def:g6:shapes:quadrilaterals}{losango} com um \omterm{def:g3:shapes:rightangle}{ângulo reto} é um quadrado.
\end{enumerate}
\end{exercise}

\begin{exercise}[$\star$]\label{exo:g6:shapes:7}
Construa um retângulo $ABCD$ com $AB = 5$ cm e $BC = 3$ cm. Meça as
duas diagonais dele: o que você observa?
\end{exercise}

\begin{exercise}[$\star$]\label{exo:g6:shapes:8}
Construa um \omterm{def:g6:shapes:quadrilaterals}{losango} com diagonais de $6$ cm e $4$ cm. (Use a
\cref{prop:g6:shapes:diagonals}: trace antes as duas diagonais ---
\omterm{def:g6:lines:perp}{perpendiculares}, cruzando-se nos pontos médios --- e depois ligue as
quatro pontas.)
\end{exercise}

\begin{exercise}[$\star\star$]\label{exo:g6:shapes:9}
Escreva um programa de construção para um \omterm{def:g6:shapes:triangles}{triângulo isósceles} de base
$4$ cm e lados iguais de $6$ cm, no estilo do
\cref{ex:g6:shapes:program}. Troque com um colega: ele consegue
segui-lo sem ver a sua figura?
\end{exercise}

\begin{exercise}[$\star\star$]\label{exo:g6:shapes:10}
Trace um círculo de centro $O$ e um dos \omterm{def:g4:geometry:circle}{diâmetros} dele, $[AB]$. Marque
um ponto $M$ qualquer sobre o círculo e trace o triângulo $AMB$. Meça o
ângulo $\widehat{AMB}$. Mude $M$ de lugar e meça de novo. O que você
conjectura? (Esse fato surpreendente é demonstrado no
\cref{ch:g8:cosine}.)
\end{exercise}

\begin{exercise}[$\star\star$]\label{exo:g6:shapes:11}
Um \omterm{def:g4:geometry:polygon}{quadrilátero} tem diagonais de mesmo comprimento que se cortam no
ponto médio comum delas. Descubra que forma ele tem de ser, usando a
\cref{prop:g6:shapes:diagonals}: retângulo, \omterm{def:g6:shapes:quadrilaterals}{losango} ou quadrado?
(Cuidado: só uma propriedade foi dada.)
\end{exercise}

\begin{exercise}[$\star\star\star$]\label{exo:g6:shapes:12}
Numa malha de $4 \times 4$ pontos, quantos quadrados diferentes você
consegue traçar com os cantos sobre pontos da malha --- contando também
os quadrados ``inclinados''? Trace ao menos um quadrado inclinado e
explique por que os lados dele são iguais.
\end{exercise}

\section{Problema: o que as diagonais sabem}

\begin{problem}[{Problema de fim de semana --- reconstruir cada
quadrilátero especial só a partir das diagonais, e contar as diagonais
de um polígono qualquer}]
\label{pb:g6:shapes:1}
A \cref{prop:g6:shapes:diagonals} lê uma forma e conta o que se passa
com as diagonais dela. Este problema joga o jogo \emph{ao contrário}:
dados apenas dois \omterm{def:g6:lines:objects}{segmentos} que se cruzam --- as futuras diagonais ---,
que forma aparece quando ligamos as quatro pontas deles? Você vai
montar um ``dicionário das diagonais'' completo, entender \emph{por
que} a construção do \omterm{def:g6:shapes:quadrilaterals}{losango} do \cref{exo:g6:shapes:8} funciona,
encontrar uma forma que ainda não tem nome e terminar contando as
diagonais de \omterm{def:g4:geometry:polygon}{polígonos} com bem mais de quatro lados.

\medskip
\noindent\textbf{Parte I --- O dicionário das diagonais.}
Em cada questão, trace os dois \omterm{def:g6:lines:objects}{segmentos} como descrito, ligue as quatro
pontas em ordem e identifique o \omterm{def:g4:geometry:polygon}{quadrilátero} obtido.
\begin{enumerate}
  \item Dois \omterm{def:g6:lines:objects}{segmentos} de $6$ cm cada, \omterm{def:g6:lines:perp}{perpendiculares}, cruzando-se no
        ponto médio comum. Meça os quatro lados e os quatro ângulos da
        forma resultante: o que é ela?
  \item Dois \omterm{def:g6:lines:objects}{segmentos} de $8$ cm e $5$ cm, \omterm{def:g6:lines:perp}{perpendiculares},
        cruzando-se no ponto médio comum. Que forma aparece
        (\cref{exo:g6:shapes:8})?
  \item Dois \omterm{def:g6:lines:objects}{segmentos} de $7$ cm cada --- iguais, mas \emph{não}
        \omterm{def:g6:lines:perp}{perpendiculares} --- cruzando-se no ponto médio comum. Meça os
        lados e os ângulos: que forma você conjectura?
  \item Dois \omterm{def:g6:lines:objects}{segmentos} de $8$ cm e $5$ cm, não \omterm{def:g6:lines:perp}{perpendiculares},
        cruzando-se no ponto médio comum. A forma obtida não é nenhuma
        das três da \cref{def:g6:shapes:quadrilaterals} --- descreva o
        que você vê (que lados parecem iguais? paralelos?). Esse
        \omterm{def:g4:geometry:polygon}{quadrilátero} recebe nome no \cref{ch:g7:central}.
  \item Resuma as questões 1 a 4 num dicionário de quatro entradas:
        diagonais \emph{iguais / \omterm{def:g6:lines:perp}{perpendiculares} / as duas coisas /
        nenhuma das duas} (sempre cruzando-se nos pontos médios)
        $\rightarrow$ nome (ou descrição) da forma.
\end{enumerate}

\medskip
\noindent\textbf{Parte II --- Por que as construções funcionam.}
\begin{enumerate}[resume]
  \item Na questão 2, as duas diagonais cortam a forma em quatro
        \omterm{def:g6:shapes:triangles}{triângulos retângulos} em torno do ponto de cruzamento. Quais
        são os dois comprimentos dos catetos de cada um desses
        triângulos? Explique por que os quatro triângulos são cópias
        exatas uns dos outros e deduza por que os quatro lados da forma
        são iguais --- o segredo do \cref{exo:g6:shapes:8}.
  \item Deduza um teste só com o compasso: depois de construir uma
        forma qualquer a partir das diagonais, como conferir apenas com
        o compasso (sem medir com régua) que os quatro lados são
        iguais?
  \item Escreva um programa de construção
        (\cref{met:g6:shapes:program}) para um quadrado cujas
        \emph{diagonais} meçam $6$ cm --- e não os lados.
  \item No quadrado da questão 8, o lado é maior ou menor que $3$ cm
        --- a \omterm{def:g3:division:half}{metade} da diagonal? Meça para descobrir. \omterm{def:g3:division:half}{Metade} da
        resposta já pode ser \emph{demonstrada}: ir de um canto ao
        seguinte passando pelo centro é um desvio de $3 + 3$ cm, e um
        caminho reto é sempre mais curto que um desvio --- então o lado
        é menor que $6$ cm. (Demonstrar que ele também supera $3$ cm, e
        calculá-lo exatamente, é tarefa do
        \cref{ch:g8:pythagoras}.)
  \item Uma entrada nova para o dicionário: trace um \omterm{def:g6:lines:objects}{segmento} $[AC]$ de
        $8$ cm; sobre ele, marque o ponto $H$ a $3$ cm de $A$; trace por
        $H$ o \omterm{def:g6:lines:objects}{segmento} perpendicular $[BD]$ de $5$ cm, com $H$ como
        ponto médio \emph{dele}. Ligue $ABCD$. Essa forma é uma
        \emph{pipa}\index{pipa} (uma diagonal corta a outra ao meio,
        formando \omterm{def:g3:shapes:rightangle}{ângulo reto}). Usando os triângulos dos cantos como na
        questão 6, descubra que lados de uma \omterm{pb:g6:shapes:1}{pipa} são iguais a quais.
\end{enumerate}

\medskip
\noindent\textbf{Parte III --- Contar diagonais.}
Uma \emph{diagonal} de um \omterm{def:g4:geometry:polygon}{polígono} liga dois \omterm{def:g5:solids:def}{vértices} que não são
vizinhos.
\begin{enumerate}[resume]
  \item Trace um \omterm{def:g4:geometry:polygon}{quadrilátero}, um \omterm{def:g4:geometry:polygon}{pentágono} (cinco \omterm{def:g5:solids:def}{vértices}) e um
        \omterm{def:g4:geometry:polygon}{hexágono} (seis \omterm{def:g5:solids:def}{vértices}) e trace todas as diagonais deles.
        Conte-as: quantas em cada caso?
  \item Contando sem desenhar, para o \omterm{def:g4:geometry:polygon}{hexágono}: de um \omterm{def:g5:solids:def}{vértice}, quantas
        diagonais saem? (A quantos dos seis \omterm{def:g5:solids:def}{vértices} ele \emph{não}
        pode ser ligado por uma diagonal?) Multiplicar pelos seis
        \omterm{def:g5:solids:def}{vértices} conta cada diagonal exatamente duas vezes --- por quê?
        Recupere a sua contagem da questão 11. (A mesma contagem em
        dobro apareceu nos apertos de mão do \cref{pb:g6:wholes:1}.)
  \item Use o método da questão 12 para contar as diagonais de um
        \omterm{def:g4:geometry:polygon}{polígono} com $10$ \omterm{def:g5:solids:def}{vértices} e depois com $20$ \omterm{def:g5:solids:def}{vértices}.
  \item Um \omterm{def:g4:geometry:polygon}{polígono} tem exatamente $44$ diagonais. Quantos \omterm{def:g5:solids:def}{vértices} ele
        tem? (Experimente o método da questão 13 em alguns
        candidatos.)
  \item Confira o seu método no menor \omterm{def:g4:geometry:polygon}{polígono} de todos: quantas
        diagonais tem um triângulo, e o que responde o método de
        contagem da questão 12? Conclua com uma frase sobre o que
        apertos de mão e diagonais têm em comum.
\end{enumerate}
\end{problem}
